\section{Consequences for literal canonical assignments}
\label{sec:source-applications}

The numerical results apply to original, undecorated G\"odel-$T$ terms.
Fix one computable admissible policy for expanding the finite sums and
repetitions in the assignment clauses.  The canonical assignment here is
the resulting zero-history assignment of the literal input term.
Normalizing the term first would define a different object.  The type
level is
\[
 \lv(o)=0,\qquad
 \lv(\sigma\to\tau)=\max\{\lv(\sigma)+1,\lv(\tau)\}.
\]
Vectors are padded by zeros when their levels differ.  At any admissible
valuation, weak comparison means weakness in every coordinate;
Wilken--Weiermann strict comparison means strictness in coordinate zero
and weakness in every positive coordinate.  Equality means equality in
every coordinate.  The strict relation is not a lexicographic order.

We use the following exact source interface.  Its domain and construction
are established from the original assignment clauses in
\cite[the canonical-source section]{wwcomparison}; the source notation
and original bounded-norm conditions are those of \cite{ww}.

\begin{proposition}[Canonical forward interface]
\label{prop:source-interface}
Let $A$ be an original well-typed term all of whose free variables have
ground type.  Its literal canonical vector is computable.  Every positive
coordinate is a closed computable ordinal, and its scalar has an
effectively constructed, valuation-preserving expression in
$\mathcal P_d$, where the $d$ inputs are the distinct ground free names.
Repeated names remain one input.  Arbitrary finite internal and bound
types, result types, recursor instances, and nesting are allowed.
Every positive integer is an admissible value of a ground component,
independently for distinct ground names.
\end{proposition}

The relevant invariant explains the interface's scope.  An actual positive
coordinate is generated from $1,\omega$, positive-index component
variables, natural sums, and graded terms
$2^a\otimes b$, where the exponent comes from the next coordinate.
It contains no promotion of a scalar collapse to a positive coordinate.
The box operation, partials and abstraction preserve this invariant
simultaneously.  Abstraction removes its bound component variables;
ground free names have no positive components.  Hence positive outputs
are closed.  An actual scalar contains only its scalar variables,
positive constants, addition, fixed finite repetitions and
\[
 \operatorname{Psi}(c,u)=\psi(\omega\otimes c+u).
\]
Its positive head $c$ is closed, and $\omega\otimes c$ is a fixed
nonzero limit.  The syntactic abstraction multiplier is fixed before
valuation.  These facts give precisely the stated positive expression
and preserve correlations.  The longer domain proof, including the
correct weak scalar partial estimate and the original abstraction
budget, is the cited source theorem; membership in an unrestricted
ambient ordinal-expression grammar would not suffice.

\subsection{Equality and one shared ground name}

\begin{corollary}[Universal source equality]
\label{cor:source-equality}
Universal equality of literal canonical vectors is decidable for every
original source pair with ground free variables, with no bound on its
internal types.
\end{corollary}
\begin{proof}
Compute the two vectors and compare their finitely many closed positive
coordinates.  If they agree, apply the complete normal form of
Theorem~\ref{u:thm:identity} to the two scalar expressions.
Different forms give the computed positive countervaluation from that
theorem.  The comparison substitutes admissible ground coordinates;
it does not substitute source numerals and recompute a different
canonical assignment.
\end{proof}

\begin{corollary}[One-name Boolean comparison]
\label{cor:source-one-name}
Take a finite Boolean combination of pointwise canonical equality, weak
comparison, and Wilken--Weiermann strict comparison atoms.  If the union
of its free names is contained in one ground name $X:o$, either its
existential or universal closure over the admissible component of $X$
is decidable.  Internal types and result levels are unrestricted.
\end{corollary}
\begin{proof}
For an atom comparing $A$ and $B$, compute its positive-coordinate guard
\[
 U_{\le}=\bigwedge_{i>0}\can A_i\le\can B_i,
 \qquad
 U_=\bigwedge_{i>0}\can A_i=\can B_i,
\]
using zero padding.  The three pointwise atoms are respectively
\[
 U_=\wedge(E_A=E_B),\qquad
 U_{\le}\wedge(E_A\le E_B),\qquad
 U_{\le}\wedge(E_A<E_B).
\]
Place each computed guard inside its original Boolean atom and apply
Corollary~\ref{u:cor:boolean}.  In particular a false guard is not
removed from beneath a negation.  The admissible ground domain is exactly
$\NN$ by Proposition~\ref{prop:source-interface}.
\end{proof}

One name in each term is not sufficient if the two names differ: the
one-input hypothesis concerns the whole query's shared input set.

\subsection{Finite coefficients and recursor-free terms}

\begin{theorem}[The finite-coefficient class]
\label{thm:source-finite}
For a ground-open source pair, suppose every literal scalar collapse
$\psi(\omega\otimes c+g)$ in the generated assignments has a finite
closed coefficient $c$.  Universal weak comparison, Wilken--Weiermann
strict comparison, and equality are decidable for the exact standard
control.  Every recursor-free ground-open source satisfies this condition,
regardless of its internal type levels.
\end{theorem}
\begin{proof}
The hypothesis is decidable by finite traversal of the actual scalar
trees and closed ordinal comparison with $\omega$.  For $c=k<\omega$,
Corollary~\ref{f:cor:finite-coefficient} replaces the call exactly by
\[
          g\longmapsto f^k(g+2k)+k.
\]
The iterate length is a closed integer determined from the source.  Its
expansion is therefore a finite expression in $(\N;0,1,+,<,f)$,
not a variable-iteration primitive.  Retain every original name and apply
Theorem~\ref{m:thm:snyatkov} to the required universal scalar formula,
with $x_i\ge1$.  The closed positive-coordinate tests complete the
corresponding vector comparison.

For the recursor-free assertion, inspect the strengthened positive
coordinate invariant in Proposition~\ref{prop:source-interface} before
erasing the bound components.  The only source constant introducing
$\omega$ into a positive coordinate is a recursor.  Without recursors,
the positive construction uses $1$, component variables, natural sums
and products, and $2^a$.  After the ground-free restriction, all bound
components have been replaced or removed by the literal partials and
all-ones substitutions.  Induction through those same operations leaves
finite positive values: sums, products, and $2^n$ preserve finiteness.
This holds at every coordinate and every internal type.  Thus every
direct scalar coefficient is finite.
\end{proof}

A recursor is not itself a certificate of failure.  For example the
canonical vector of $R^oX$ is $(f(x+4)+1,1,\omega)$; its displayed
scalar collapse has coefficient one.  The theorem checks the actual
coefficient condition rather than a syntactic ban on all recursors.

\subsection{Polynomial heads and internal type level two}

\begin{lemma}[Polynomial direct heads]
\label{lem:source-polynomial}
If every internal type and subtype in a ground-open source has level at
most two, all its direct scalar heads are polynomial limits below
$\omega^\omega$.
\end{lemma}
\begin{proof}
Follow the same generated positive-coordinate invariant.  There is no
coordinate three available.  At coordinate two an exponent constructor
would require such a coordinate, so its closed values are natural linear
combinations of $1$ and $\omega$.  The positive box and abstraction
clauses only add such values and take fixed natural multiples there.
At coordinate one the exponential factors therefore have the form
\[
                  2^{a\omega+b}=\omega^a2^b,
                  \qquad a,b<\omega.
\]
Natural sums and products of these factors are polynomials in $\omega$.
Whole-free partial copies and fixed repetitions preserve this property;
temporary component variables disappear in abstraction.  It follows
that every closed coefficient $c$ of a direct scalar collapse is a
polynomial in $\omega$, so $\omega\otimes c$ is a polynomial limit
below $\omega^\omega$.

The hypothesis includes the full types of recursor constants and of
temporary binders, not just the result type or the free context.  A
higher-level recursor could introduce a coordinate above two and would
invalidate this induction even if the final result were ground.
\end{proof}

\begin{theorem}[Comparison through internal level two]
\label{thm:source-leveltwo}
Universal weak comparison, Wilken--Weiermann strict comparison, and
equality of literal canonical assignments are decidable for ground-open
sources whose maximum internal type level is at most two.
\end{theorem}
\begin{proof}
Compute the closed positive-coordinate tests and the scalar circuit from
Proposition~\ref{prop:source-interface}.  Its heads are polynomial by
Lemma~\ref{lem:source-polynomial}.  Translate all circuit nodes by the
functional graphs in Theorem~\ref{m:thm:graphs}.  Quantify the fresh
natural graph variables together with the original positive inputs and
conjoin the graph equations positively.  The counterexample output test
is $E_A>E_B$ for weak comparison, $E_A\ge E_B$ for strict comparison,
and $E_A\ne E_B$ for equality.  Thus no graph equation is negated when
the requested relation is negated.

Lemma~\ref{m:lem:graph-domains} splits the finite graph exceptions and
puts each remaining case on its individual word domains.
Theorem~\ref{m:thm:existential} decides the resulting closed existential
sentences.  A counterexample exists exactly when the corresponding
universal scalar relation fails.  The needed degree bound is finite and
computed from the input pair; there is no fixed degree bound across all
terms of internal level two.
\end{proof}

At arbitrary internal levels the common-head, balanced, or single-core
theorems can still apply to a scalar residual carrying the corresponding
finite certificate.  One first computes the literal assignment and then
checks the certificate.  Failure of a proposed certificate is not a
negative answer to comparison.  No reverse source-realization theorem
for every auxiliary numerical expression is needed in any of these
applications.
